\documentclass[11pt,oneside]{article}

\pdfminorversion=7
\usepackage{QuizStyle}

\hypersetup{
    pdftitle={2026Fall Quiz for MATH230 Probability \& Statistics},
    pdfauthor={Jungwoo Kim},
    pdfsubject={Quiz for MATH230 Probability \& Statistics},
    pdfcreator={LaTeX}
}

\begin{document}

\quiztitle{2026Fall Quiz}

\begin{problem}{1.1}
Four married couples have bought 8 seats in the same row for a concert. In how many different ways can they be seated
\begin{problemsubparts}
    \item with no restrictions?
    \item if each couple is to sit together?
    \item if at least one couple cannot sit next to each other?
    \item if all the men sit together to the right of all the women?
\end{problemsubparts}

\textit{Note.} You can write your answer in calculator-ready form. You don't have to calculate the exact value of fractions or combinations. For example, $8!$, $(4!)^4$, and $\binom{8}{3}/\binom{10}{3}$ are acceptable forms.
\end{problem}

\begin{solution}
\begin{problemsubparts}
    \item With no restrictions, the 8 people can be arranged in
    \[
        \boxed{8!}
    \]
    ways.

    \item Treat each couple as a single block. The four couple-blocks can be arranged in $4!$ ways, and the two people within each couple can be arranged in $2!$ ways. Hence the number of arrangements is
    \[
        \boxed{4!(2!)^4}.
    \]

    \item The complement consists of the arrangements in which all four couples sit together. By parts (a) and (b), the required number is
    \[
        \boxed{8!-4!(2!)^4}.
    \]

    \item The four women must occupy the four leftmost seats and the four men the four rightmost seats. Each group can be arranged independently in $4!$ ways, giving
    \[
        \boxed{4!\cdot4!}.
    \]
\end{problemsubparts}
\end{solution}

\Needspace{21\baselineskip}
\begin{problem}{1.2}
A paint-store chain produces and sells latex and semigloss paint. Based on long-range sales, the probability that a customer will purchase latex paint is $0.25$. Of those that purchase latex paint, $75\%$ also purchase rollers. But only $25\%$ of semigloss paint buyers purchase rollers. A randomly selected buyer purchases a roller and a can of paint. What is the probability that the paint is latex?

\textit{Note.} If you introduce symbols (e.g., $A$, $B$) in your solution, you must clearly define what each symbol represents.
\end{problem}

\begin{solution}
Let
\[
    L=\{\text{the buyer purchases latex paint}\},\qquad
    S=\{\text{the buyer purchases semigloss paint}\},
\]
and
\[
    R=\{\text{the buyer purchases a roller}\}.
\]
Then
\[
    \Prob(L)=\frac14,
    \qquad
    \Prob(R\mid L)=\frac34,
    \qquad
    \Prob(R\mid S)=\frac14,
\]
and, since the paint is either latex or semigloss,
\[
    \Prob(S)=1-\Prob(L)=\frac34.
\]
By the law of total probability,
\[
\begin{aligned}
    \Prob(R)
    &=\Prob(R\mid L)\Prob(L)+\Prob(R\mid S)\Prob(S)\\
    &=\frac34\cdot\frac14+\frac14\cdot\frac34
      =\frac38.
\end{aligned}
\]
Therefore, by Bayes' rule,
\[
    \boxed{
    \Prob(L\mid R)
    =\frac{\Prob(R\mid L)\Prob(L)}{\Prob(R)}
    =\frac{(3/4)(1/4)}{3/8}
    =\frac12
    }.
\]
Thus the probability that the paint is latex is $50\%$.
\end{solution}

\Needspace{23\baselineskip}
\begin{problem}{2.1}
For a laboratory assignment, if the equipment is working, the density function of the observed outcome $X$ is
\[
    f_X(x)=
    \begin{cases}
        c(2-x), & 0<x<2,\\
        0, & \text{elsewhere}.
    \end{cases}
\]
\begin{problemsubparts}
    \item Find the value of $c$ so that $f_X$ is a valid density function.
    \item Find $\Prob(X<1)$.
    \item Given that $X\geq1$, what is the probability that $X$ will be less than $1.5$?
\end{problemsubparts}
\end{problem}

\begin{solution}
\begin{problemsubparts}
    \item A probability density must integrate to $1$. Thus
    \[
        1=\int_{-\infty}^{\infty}f_X(x)\,\dd x
        =\int_0^2 c(2-x)\,\dd x
        =2c,
    \]
    so
    \[
        \boxed{c=\frac12}.
    \]

    \item Using $c=1/2$,
    \[
        \boxed{
        \Prob(X<1)
        =\int_0^1 \frac12(2-x)\,\dd x
        =\frac34
        }.
    \]

    \item By the definition of conditional probability,
    \[
        \Prob(X<1.5\mid X\geq1)
        =\frac{\Prob(1\leq X<1.5)}{\Prob(X\geq1)}.
    \]
    Here,
    \[
        \Prob(1\leq X<1.5)
        =\int_1^{3/2}\frac12(2-x)\,\dd x
        =\frac3{16},
    \]
    while
    \[
        \Prob(X\geq1)=1-\Prob(X<1)=\frac14.
    \]
    Therefore,
    \[
        \boxed{\Prob(X<1.5\mid X\geq1)=\frac{3/16}{1/4}=\frac34}.
    \]
\end{problemsubparts}
\end{solution}

\Needspace{29\baselineskip}
\begin{problem}{2.2}
The amount of kerosene, in thousands of liters, in a tank at the beginning of any day is a random amount $Y$ from which a random amount $X$ is sold during that day. Suppose that the tank is not resupplied during the day so that $x\leq y$, and assume that the joint density function of these variables is
\[
    f_{X,Y}(x,y)=
    \begin{cases}
        2, & 0<x\leq y<1,\\
        0, & \text{elsewhere}.
    \end{cases}
\]
\begin{problemsubparts}
    \item Find the marginal density function of $X$.
    \item Find the marginal density function of $Y$.
    \item Determine if $X$ and $Y$ are independent.
    \item Find $\Prob(1/2<Y<3/4\mid X=1/4)$.
\end{problemsubparts}
\end{problem}

\begin{solution}
\begin{problemsubparts}
    \item For fixed $x$ with $0<x<1$, the support has $x\leq y<1$. Hence
    \[
        f_X(x)=\int_x^1 f_{X,Y}(x,y)\,\dd y
        =\int_x^1 2\,\dd y
        =2(1-x).
    \]
    Therefore,
    \[
        \boxed{
        f_X(x)=
        \begin{cases}
            2(1-x), & 0<x<1,\\
            0, & \text{elsewhere}.
        \end{cases}}
    \]

    \item For fixed $y$ with $0<y<1$, the support has $0<x\leq y$. Thus
    \[
        f_Y(y)=\int_0^y f_{X,Y}(x,y)\,\dd x
        =\int_0^y 2\,\dd x
        =2y.
    \]
    Therefore,
    \[
        \boxed{
        f_Y(y)=
        \begin{cases}
            2y, & 0<y<1,\\
            0, & \text{elsewhere}.
        \end{cases}}
    \]

    \item If $X$ and $Y$ were independent, then $f_{X,Y}(x,y)=f_X(x)f_Y(y)$ almost everywhere. But for $0<y<x<1$,
    \[
        f_{X,Y}(x,y)=0,
        \qquad
        f_X(x)f_Y(y)=4y(1-x)>0.
    \]
    Hence
    \[
        \boxed{X\text{ and }Y\text{ are not independent}.}
    \]

    \item For $0<x<1$, the conditional density of $Y$ given $X=x$ is
    \[
        f_{Y\mid X}(y\mid x)
        =\frac{f_{X,Y}(x,y)}{f_X(x)}
        =\frac{1}{1-x},
        \qquad x<y<1.
    \]
    When $X=1/4$,
    \[
        f_{Y\mid X}(y\mid1/4)=\frac43,
        \qquad \frac14<y<1.
    \]
    Therefore,
    \[
        \boxed{
        \Prob\left(\frac12<Y<\frac34\,\middle|\,X=\frac14\right)
        =\int_{1/2}^{3/4}\frac43\,\dd y
        =\frac13
        }.
    \]
\end{problemsubparts}
\end{solution}

\Needspace{16\baselineskip}
\begin{problem}{3.1}
The probability distribution of the discrete random variable $X$ is
\[
    p_X(x)=\binom{5}{x}
    \left(\frac15\right)^x
    \left(\frac45\right)^{5-x},
    \qquad x=0,1,2,3,4,5.
\]
Find the mean of $X$.
\end{problem}

\begin{solution}
By the definition of expectation for a discrete random variable,
\[
\begin{aligned}
    \Exp[X]
    &=\sum_{x=0}^{5}x\,p_X(x)\\
    &=\sum_{x=0}^{5}
      x\binom{5}{x}
      \left(\frac15\right)^x
      \left(\frac45\right)^{5-x}\\
    &=\frac{1280+1280+480+80+5}{3125}\\
    &=\boxed{1}.
\end{aligned}
\]
\end{solution}

\Needspace{20\baselineskip}
\begin{problem}{3.2}
Let $X$ and $Y$ be random variables with joint density function
\[
    f_{X,Y}(x,y)=
    \begin{cases}
        4xy, & 0<x,y<1,\\
        0, & \text{elsewhere}.
    \end{cases}
\]
\begin{problemsubparts}
    \item Find the expected value of $X+Y$.
    \item Find the expected value of $XY$.
\end{problemsubparts}
\end{problem}

\begin{solution}
\begin{problemsubparts}
    \item By the definition of expectation for a function of jointly continuous random variables,
    \[
    \begin{aligned}
        \Exp[X+Y]
        &=\int_0^1\int_0^1 (x+y)f_{X,Y}(x,y)\,\dd y\,\dd x\\
        &=\int_0^1\int_0^1 (4x^2y+4xy^2)\,\dd y\,\dd x\\
        &=\int_0^1\left(2x^2+\frac43x\right)\,\dd x\\
        &=\boxed{\frac43}.
    \end{aligned}
    \]

    \item Again using the definition of expectation,
    \[
    \begin{aligned}
        \Exp[XY]
        &=\int_0^1\int_0^1 xy\,f_{X,Y}(x,y)\,\dd y\,\dd x\\
        &=4\int_0^1\int_0^1 x^2y^2\,\dd y\,\dd x\\
        &=4\left(\int_0^1x^2\,\dd x\right)
             \left(\int_0^1y^2\,\dd y\right)\\
        &=\boxed{\frac49}.
    \end{aligned}
    \]
\end{problemsubparts}
\end{solution}

\Needspace{24\baselineskip}
\begin{problem}{4.1}
From a sack of fruit containing 2 oranges, 1 apple, and 2 bananas, a random sample of 3 pieces of fruit is selected. If $X$ is the number of oranges and $Y$ is the number of apples in the sample, find the covariance of the random variables $X$ and $Y$.
\end{problem}

\begin{solution}
There are $\binom{5}{3}=10$ equally likely samples. To obtain $X=x$ oranges and $Y=y$ apples, the sample must contain $3-x-y$ bananas. Hence the joint probability distribution is
\[
    p_{X,Y}(x,y)
    =\frac{\binom{2}{x}\binom{1}{y}\binom{2}{3-x-y}}{\binom{5}{3}}
\]
for the admissible values of $x$ and $y$, giving
\[
\begin{array}{c@{\qquad}ccc}
\toprule
p_{X,Y}(x,y)&\multicolumn{3}{c}{x}\\
\cmidrule(lr){2-4}
y&0&1&2\\
\midrule
0&0&\frac{2}{10}&\frac{2}{10}\\
1&\frac{1}{10}&\frac{4}{10}&\frac{1}{10}\\
\bottomrule
\end{array}
\]
The marginal probabilities give
\[
\begin{aligned}
    \Exp[X]
    &=(1)\left(\frac{2+4}{10}\right)
      +(2)\left(\frac{2+1}{10}\right)
      =\frac65,\\
    \Exp[Y]
    &=(1)\left(\frac{1+4+1}{10}\right)
      =\frac35.
\end{aligned}
\]
Also, only the pairs $(1,1)$ and $(2,1)$ contribute to $\Exp[XY]$:
\[
\begin{aligned}
    \Exp[XY]
    &=\sum_x\sum_y xy\,p_{X,Y}(x,y)\\
    &=(1)(1)\frac4{10}+(2)(1)\frac1{10}
      =\frac35.
\end{aligned}
\]
Therefore, by the covariance formula,
\[
    \boxed{
    \Cov(X,Y)
    =\Exp[XY]-\Exp[X]\Exp[Y]
    =\frac35-\left(\frac65\right)\left(\frac35\right)
    =-\frac3{25}
    }.
\]
\end{solution}

\Needspace{23\baselineskip}
\begin{problem}{4.2}
Let $X$ and $Y$ be random variables with joint density function
\[
    f_{X,Y}(x,y)=
    \begin{cases}
        \dfrac12, & 0<x\leq y<2,\\
        0, & \text{elsewhere}.
    \end{cases}
\]
\begin{problemsubparts}
    \item Find $\Var(X)$ and $\Var(Y)$.
    \item Find $\Cov(X,Y)$.
    \item Find $\rho_{X,Y}$, the correlation coefficient between $X$ and $Y$.
    \item Find $\Var(X+Y)$.
\end{problemsubparts}

\textit{Hint:} You may use $\Exp[X]=\frac23$, $\Exp[Y]=\frac43$, $\Exp[X^2]=\frac23$, and $\Exp[Y^2]=2$.
\end{problem}

\begin{solution}
The support is the triangular region $0<x\leq y<2$.
\begin{problemsubparts}
    \item Using the given moments and the variance formula,
    \[
    \begin{aligned}
        \Var(X)
        &=\Exp[X^2]-\bigl(\Exp[X]\bigr)^2
          =\frac23-\left(\frac23\right)^2
          =\boxed{\frac29},\\
        \Var(Y)
        &=\Exp[Y^2]-\bigl(\Exp[Y]\bigr)^2
          =2-\left(\frac43\right)^2
          =\boxed{\frac29}.
    \end{aligned}
    \]

    \item We first calculate the mixed moment using the joint density:
    \[
    \begin{aligned}
        \Exp[XY]
        &=\int_0^2\int_0^y xy\,f_{X,Y}(x,y)\,\dd x\,\dd y\\
        &=\int_0^2\int_0^y\frac{xy}{2}\,\dd x\,\dd y\\
        &=\frac14\int_0^2y^3\,\dd y
          =1.
    \end{aligned}
    \]
    Therefore,
    \[
        \boxed{\Cov(X,Y)
        =\Exp[XY]-\Exp[X]\Exp[Y]
        =1-\left(\frac23\right)\left(\frac43\right)
        =\frac19}.
    \]

    \item Since both variances are positive, the correlation coefficient is
    \[
        \boxed{\rho_{X,Y}
        =\frac{\Cov(X,Y)}{\sqrt{\Var(X)\Var(Y)}}
        =\frac{1/9}{\sqrt{(2/9)(2/9)}}
        =\frac12}.
    \]

    \item Using the variance of a sum,
    \[
    \begin{aligned}
        \Var(X+Y)
        &=\Var(X)+\Var(Y)+2\Cov(X,Y)\\
        &=\frac29+\frac29+2\left(\frac19\right)
          =\boxed{\frac23}.
    \end{aligned}
    \]
\end{problemsubparts}
\end{solution}

\end{document}
